\subsection{A stabilizer obstruction to scalar pivotal rotation}
\label{sec:stabilizer-rotation}

\begin{proposition}[Anticommuting stabilizers and third rotation]
\label{prop:stabilizer-rotation}
Let $\mathcal C$ be a spherical fusion category, let $g$ be an invertible
of order two for which $A=\mathbf1\oplus g$ is an algebra, and let $Y$
be a simple object with $gY\cong Y\cong Yg$. Suppose that left
multiplication by $g$ exchanges the two right $A$-module structures
on $Y$. Put $m=\dim\operatorname{Hom}(\mathbf1,Y^{\otimes3})$.
Then $m$ is even. If $m>0$, the pivotal third-rotation operator on this
space is not scalar. More precisely,
\[
 \nu_3(Y)=-(a+b\zeta_3+c\zeta_3^2),
 \qquad a,b,c\in\ZZ_{\ge0},\quad a+b+c=m/2.
\]
If $3$ is coprime to the Frobenius--Schur exponent, this becomes
\[
 \nu_3(Y)\in
 \{-m/2+3b:0\le b\le\lfloor m/4\rfloor\}.
\]
In particular, if $m=2$ under this coprimality hypothesis, then
\[
 \nu_3(Y)=-1.
\]
\end{proposition}
\begin{proof}
Use strict monoidal notation and an associative identification
$q:g^2\xrightarrow\sim\mathbf1$ supplied by the algebra $A$.
Choose isomorphisms $l:gY\to Y$ and $r:Yg\to Y$, normalized by
\[
 l(\operatorname{id}_g\otimes l)=q\otimes\operatorname{id}_Y,
 \qquad r(r\otimes\operatorname{id}_g)=\operatorname{id}_Y\otimes q.
\]
The two normalized right module structures are $r$ and $-r$.
Transporting the right action on $gY$ through $l$ gives
\[
 r'=l(\operatorname{id}_g\otimes r)(l^{-1}\otimes\operatorname{id}_g).
\]
The hypothesis says $r'=-r$, or equivalently
\begin{equation}\label{eq:stabilizer-anticommutator}
 r(l\otimes\operatorname{id}_g)=-l(\operatorname{id}_g\otimes r).
\end{equation}
This also shows directly how the hypothesis is read from the
anti-invariant module matrices: it is $T_g[Y,Y]=-1$.

Define an invertible endomorphism of $Y^2$ by
\[
 a=(r\otimes l)(\operatorname{id}_Y\otimes q^{-1}\otimes\operatorname{id}_Y),
\]
and put $P=a\otimes\operatorname{id}_Y$, $Q=\operatorname{id}_Y\otimes a$ on $Y^3$.
Expanding either composite inserts two copies of $q^{-1}$, with common
intermediate tensor word $YggYggY$. The endpoint maps agree. On the
middle $gYg$, the two composites use the two sides of
\eqref{eq:stabilizer-anticommutator}. Consequently
\[
 PQ=-QP.
\]
The restrictions to $V=\operatorname{Hom}(\mathbf1,Y^3)$ remain
invertible. Taking determinants shows that $\dim V$ is even.

Let $E$ be pivotal cyclic rotation moving the last $Y$ to the first
position. The canonical cyclic map
\[
 \operatorname{Hom}(\mathbf1,B\otimes Y)
 \longrightarrow \operatorname{Hom}(\mathbf1,Y\otimes B)
\]
is natural in the entire object $B$. Applying this naturality to
$B=Y^2$ and $a\in\operatorname{End}(B)$ gives the exact identity
\[
 EP=QE.
\]
The construction and its naturality are those of
\cite[Definition~3.1, Lemma~3.4 and Theorem~5.1]{NSpiv}.
No braiding or self-duality sign is used. If $E$ were scalar on nonzero
$V$, then $P=Q$, contradicting $PQ=-QP$ and invertibility.

The module-action normalizations also give $a^2=I$: after two pair
insertions, the first endpoint applies $r(r\otimes\operatorname{id}_g)$,
the second applies $l(\operatorname{id}_g\otimes l)$, and the two
resulting copies of $q$ cancel the inserted copies of $q^{-1}$ by
associativity. Set
$A_j=E^jPE^{-j}$ for $j\in\ZZ/3\ZZ$. These are involutions, and the
anticommutation of $A_0=P$ and $A_1=Q$, together with cyclic conjugation,
makes them pairwise anticommuting. They therefore give a representation
of the complex Clifford algebra $\mathrm{Cl}_3\cong M_2(\CC)\oplus
M_2(\CC)$. In that algebra the element
\[
 U=\frac{-I+A_0A_1+A_1A_2+A_2A_0}{2}
\]
satisfies $U^3=I$, $UA_jU^{-1}=A_{j+1}$, and has trace $-1$ on each
two-dimensional simple module. These identities also follow by direct
multiplication using $A_i^2=I$ and $A_iA_j=-A_jA_i$.
The eight-dimensional universal algebra calculation is reproduced with
rational coefficients by
\filename{research/round2_audit/clifford_rotation.py}; it verifies
$U^3=I$ and $UA_j=A_{j+1}U$. Its scalar coefficient is $-1/2$, while
the three bivectors are traceless on either two-dimensional simple,
giving trace $-1$ without a matrix-choice assumption.
Thus $S=U^{-1}E$ commutes with the Clifford action and satisfies
$S^3=I$. On its multiplicity spaces, of total dimension $m/2$, let
$a,b,c$ be the multiplicities of $1,\zeta_3,\zeta_3^2$ for $S$.
Taking the trace of $E=US$ proves the asserted formula. Reversing the
cyclic convention exchanges $b$ and $c$, so the same value set applies
to the standard indicator convention.

On the other hand, the central-twist
formula puts $\nu_3(Y)$ in $\QQ(\zeta_N)$, where $N$ is the
Frobenius--Schur exponent. If $3\nmid N$, then
$\QQ(\zeta_N)\cap\QQ(\zeta_3)=\QQ$. Hence $b=c$, and
$a+2b=m/2$ gives the rational-value list.
\end{proof}

\begin{remark}[No self-duality requirement]
The stabilizer-rotation proposition applies equally to a non-self-dual
simple $Y$. Its hypotheses and proof use isomorphisms $gY\cong Y\cong Yg$,
not an isomorphism $Y\cong Y^*$. In particular, rigidity gives
\[
 m=\dim\operatorname{Hom}(\mathbf1,Y^{\otimes3})=N_{Y,Y}^{Y^*};
\]
one must not replace the last index by $Y$ unless $Y$ is self-dual.
The pair-insertion operators act on $Y^3$ itself, and naturality of pivotal
rotation applies to $Y^2$ and its endomorphism $a$ without dualizing $Y$.
Thus the same Clifford multiplicity formula applies to each member of a
dual pair separately.
\end{remark}

\begin{corollary}[Projected third-indicator spectrum]
\label{cor:stabilizer-average}
Under the hypotheses of Proposition~\ref{prop:stabilizer-rotation},
\[
 -m/2\le\operatorname{Re}\nu_3(Y)\le m/4,
 \qquad
 \operatorname{Re}\nu_3(Y)\in
 \{m/4-3a/2: a=0,1,\ldots,m/2\}.
\]
Let $K$ be any number field satisfying
$K\cap\mathbb Q(\zeta_3)=\mathbb Q$. The normalized relative trace of
$\nu_3(Y)$ to $K$ equals its real part and therefore belongs to the
displayed finite set. This includes every real number field and the
complex cyclotomic field $\mathbb Q(\zeta_8)$.
If instead $K$ contains $\zeta_3$, the unprojected possibilities are the
full set
\[
 \{-(a+b\zeta_3+c\zeta_3^2):
       a,b,c\in\mathbb Z_{\ge0},\ a+b+c=m/2\}.
\]
\end{corollary}
\begin{proof}
The preceding proposition gives
$\nu_3(Y)=-(a+b\zeta_3+c\zeta_3^2)$ with $a+b+c=m/2$.
Its real part is $-a+(b+c)/2=m/4-3a/2$, which gives the stated set and
bounds. Under the intersection hypothesis, $K(\zeta_3)/K$ is quadratic
and its nontrivial automorphism exchanges $\zeta_3$ and $\zeta_3^2$.
Consequently its normalized relative trace sends both roots to $-1/2$,
and sends the indicator to $-a+(b+c)/2$. Transitivity gives the same
value when the trace is taken from any larger common field.
The quadratic field $\mathbb Q(\zeta_3)=\mathbb Q(\sqrt{-3})$ does not
embed into $\mathbb Q(\zeta_8)$: the latter's three quadratic subfields
are $\mathbb Q(i)$, $\mathbb Q(\sqrt2)$ and
$\mathbb Q(\sqrt{-2})$. This proves the stated complex-field example.
When $K$ contains $\zeta_3$, normalized trace fixes the entire indicator,
so the original Clifford multiplicity set must be retained.
\end{proof}


\begin{corollary}\label{cor:frontier-third-indicators}
Every pseudo-unitary categorification of $\RR{3}{215}$, in the basis
of dimensions $(1,1,d,d,d-1,d+1)$, must satisfy
\[
 \nu_3(b_4)=\nu_3(b_5)=-1.
\]
\end{corollary}
\begin{proof}
The exact full Morita-action calculation for $A=1+g$ leaves two branches,
both with $T_g=-I$ on the two fixed simples $b_4,b_5$.
They satisfy $N_{44}^{4}=N_{55}^{5}=2$, and both are self-dual.
Their third-rotation spaces therefore have dimension two.
The global dimension is $136+32\sqrt{17}$, whose norm is
$1088=2^6\cdot17$. The Cauchy theorem
\cite[Theorem~3.9]{BNRW} excludes the prime $3$ from the
Frobenius--Schur exponent. Apply
Proposition~\ref{prop:stabilizer-rotation}.
\end{proof}

\subsection{Exact certificate and exhaustive coverage}
\label{subsec:stabilizer-rotation-certificate}
The obstruction is implemented without a distinguished census label in
\filename{software/stabilizer_rotation.py}. Its input is the multiplication
tensor.  The certificate first reconstructs the complete list of
anti-invariant actions for the forced algebra $\mathbf1\oplus g$.
For the multiplicity-three frontier this list has sixteen entries.  Twelve
are rejected by exact separating functionals for the full Morita-action Gram
system, and two by its rank bound.  In each of the remaining two entries,
left multiplication by $g$ exchanges both pairs of module structures on the
fixed simples $b_4,b_5$.  Equivalently, the diagonal entries of $T_g$ on their
anti-invariant lines are both $-1$.  Their cubic invariant spaces have
 dimension two.  The pair-insertion theorem therefore forces
$\nu_3(b_4)=\nu_3(b_5)=-1$, since the Cauchy prime support is $\{2,17\}$.

The central calculation is exhaustive; it does not assume the rank-twenty-two
modular datum retained elsewhere as a relaxation witness.  All ninety-one
low induction frames are reconstructed from the tensor and its formal
codegrees.  Thirty-one frames admit character labellings; there are fifty
labellings in total and four possible second-indicator assignments.  The
certificate contains an exact Farkas vector for each of these two hundred
systems.  Each system includes the induction Gram equations, relative twist
traces through order three, cyclic-rotation character constraints, joint
Galois compatibility of dimensions and eigenvalues, and the low-row modular
relations.  It additionally imposes the two forced third indicators.  If
$Ax=b$ is such a system, its recorded integral vector $y$ is checked to satisfy
\[
 A^{\mathsf T}y\geq0,\qquad b^{\mathsf T}y<0.
\]
Every feasible central decomposition supplies $x\geq0$, so these inequalities
exclude every branch.  The variables allow all roots of unity whose prime
support is $\{2,17\}$ via their finite joint relative-trace types; no finite
upper bound on the actual exponent or central rank is imposed.  For clarity, the order-three relative-trace domain has exactly ten joint
patterns.  Write a root's order as $2^a17^b$.  If $b\geq2$ or $a\geq3$,
its normalized traces at powers $1,2,3$ all vanish.  Otherwise
$a\in\{0,1,2\}$ and $b\in\{0,1\}$ give the remaining nine distinct
patterns, with the two quadratic signs when $17$ divides the order.
For primitive seventeenth roots the first trace is
$(-1\pm\sqrt{17})/16$; the quadratic character of $2$ is $+1$ and
that of $3$ is $-1$.  This elementary enumeration explicitly includes
arbitrarily large exponent orders through the all-zero pattern.
Integral high-column decompositions are not needed for this final exclusion.

The general-scope wrapper is
\filename{software/galois_pseudo_reduction.py}. It reconstructs the full
unrestricted signed-sphericalization enumeration, including all $1480$
signed quotient tables and all sign orbits.  Every nongauge character
component is rejected by its exact arithmetic certificate, and the only
nonzero dimension-character orbit of the original ring is checked directly
to be its Frobenius--Perron orbit.  Consequently every hypothetical complex
categorification has a pseudo-unitary Galois conjugate to which the preceding
obstruction applies.  No positivity normalization is imposed during the
signed enumeration.  The combined certificate is stored in
\filename{certificates/stabilizer_rotation_215.json}; the public replay entry
point is \texttt{stabilizer\_rotation.verify\_general(R,c)}.

\begin{theorem}[Resolution of the multiplicity-three frontier]
\label{thm:frontier-noncategorifiable}\label{thm:stabilizer-rotation-frontier}
The based ring $\RR{3}{215}$ has no complex fusion categorification.
\end{theorem}
\begin{proof}
The unrestricted signed-sphericalization calculation of
Proposition~\ref{prop:frontier-galois-pu}, independently reconstructed by
the general certificate above, shows that any hypothetical
categorification has a pseudo-unitary Galois conjugate. Work with that
conjugate.

The complete involution-module domain has sixteen branches. The full
Morita-action Gram identities exclude fourteen, and either remaining
branch satisfies the hypotheses of
Corollary~\ref{cor:frontier-third-indicators}. Consequently both
$\nu_3(b_4)$ and $\nu_3(b_5)$ equal $-1$.

The exhaustive central-induction calculation in
Section~\ref{subsec:stabilizer-rotation-certificate} excludes those two
equalities for every permitted low frame, character assignment and
second-indicator choice. Each of its $200$ integral separating vectors
is checked against the reconstructed rational system; no numerical
infeasibility output is used as a certificate. Thus no pseudo-unitary
categorification exists, contradicting the reduction.
\end{proof}

This proof uses no identification of a hypothetical Drinfeld center with
a known braided category. A cyclic-eight modular datum satisfies
the central-induction conditions, but the local stabilizer and rotation
constraints prevent it,
and every other central-induction branch, from arising here.

